% Sean T. Hayes

\documentclass[14pt,aspectratio=1610]{beamer}
% \documentclass[14pt,aspectratio=1610,handout]{beamer}
\usepackage{csu-theme}
\usepackage{linked-lists}
\usepackage{arrays}
\usepackage{hhline}
\usetikzlibrary{backgrounds,calc}

\title{\Huge{}{Memory Management}\\\Large{}Static, Stack, Heap, and Pointers}
\subtitle{CSCI 315: Data Structure Analysis\\Chapter 12}
\author{Sean T. Hayes, Ph.D.}
\institute{Charleston Southern University}

%%% Comment out date to use default (today)
% \date{January 21, 2025}
\subject{Software Programming}
\keywords{C++, Static, Stack, Heap, and Pointers}

\graphicspath{{../imgs/memory-management}}

%\usepackage{tikz}
%\usetikzlibrary{arrows.meta}
%\usepackage{multirow}
\usepackage{arydshln} % for dashed lines in table

\begin{document}

\begin{frame}
  \titlepage{}
\end{frame}

\section{References \& Aliases}
\begin{frame}[fragile]{A Review of C++ References \& Aliases}
\begin{enumerate}
\item
  In a parameter declaration, \lstinline{&} makes a \textit{reference}.\\
  The parameter refers to the memory location of the original variable.

\begin{lstlisting}
void limit(double& value, double max) {
  if (value > max) {
    value = max;
  }
}
\end{lstlisting}

\begin{lstlisting}
double score = 120.5;
limit(score, 100.0);
std::cout << score << '\n'; // 100.
\end{lstlisting}
\end{enumerate}
\end{frame}

\begin{frame}[fragile]{A Review of C++ References \& Aliases}
\begin{enumerate}[<+->]
  \setcounter{enumi}{1}

\item
  In a variable declaration, \lstinline{&} makes that variable an \textit{alias}.\\
  \begin{lstlisting}
double score = 205.3;
double &alias = score;
\end{lstlisting}
  The variable is a new name for the old variable location.

\item
  Before an existing variable, \lstinline{&} evaluates to the variable's memory address.\\
  \begin{lstlisting}
double score = 91.3;
std::cout << &score; // score's address.
\end{lstlisting}
\end{enumerate}
\end{frame}

\section{Static Data}

\begin{frame}{The 3 Memory Sections (for Data Storage)}
\begin{columns}
\column{0.67\textwidth}
  \begin{itemize}[<+->]
  \item
    \textbf{Static}: storage requirements are known and allocated before execution and remain for the entire program execution.
  \item
    \textbf{Call Stack} (or execution stack): memory associated with active functions.
    \begin{itemize}
      \item Structured as \emph{stack frames} (i.e., activation records)
    \end{itemize}
  \item
    \textbf{Heap}: dynamically allocated storage; the least organized and most dynamic storage area.
  \end{itemize}
\column{0.33\textwidth}
  % % A table representing the layout of runtime memory.
% % Requires \usepackage{arydshln} % for dashed lines in the table
% % Requires \usepackage{hhline} for horizontal lines that show up with \columncolor
% \begin{tabular}{r |>{\columncolor{CSUslate}}c|}
% % \begin{tabular}{r |c|}
%   %\cline{2-2}
%   \hhline{~|-|}
%   $0$ & \textbf{Static}\\
%   $1$ & \\
%       & \\\hhline{~-}
%       & \textbf{Stack}\\
%   $a - 1$ & \\\hhline{~-}
%   \multicolumn{1}{r:}{$a$} &
%     \multicolumn{1}{>{\columncolor{CSUslate}}c:}{$\Downarrow$}\\
%     \multicolumn{1}{r:}{}  & \multicolumn{1}{>{\columncolor{CSUslate}}r:}{} \\\hhline{~-}
%   $h$ & \textbf{Heap}\\
%       &  \\
%   $n$ & \\\hhline{~-}
% \end{tabular}

\begin{tikzpicture}[memory_type/.style={minimum width=2.5cm, minimum height=1.25cm, fill=CSUcyan, text=black}, line/.style={dashed, very thick}, address/.style={font=\small, inner ysep=1pt}]
  % % Define colors for sections
  % \definecolor{staticcolor}{RGB}{173, 216, 230} % Light blue
  % \definecolor{stackcolor}{RGB}{255, 182, 193}  % Light pink
  % \definecolor{heapcolor}{RGB}{144, 238, 144}   % Light green

  % Static Memory Section
  \node[memory_type, fill=CSUcyan!50!white] (static) at (0, 0) {Static};

  % Heap Node (above Static)
  \node[memory_type, anchor=south] (heap) at (static.north) {Heap};

  % Free Node (above Heap)
  \node[memory_type, fill=none, text=white, anchor=south, align=center, font={\small\itshape}, minimum height=1.5cm] (Free) at (heap.north) {unallocated};

  % Stack Node (above Heap)
  \node[memory_type, anchor=south, align=center] (stack) at (Free.north) {Stack};

  \node[fit=(static) (stack), draw, inner sep=0pt, very thick] (Memory) {};

  % % Dashed lines where stack and heap grow.
  % \draw[line] (heap.north-|heap.west)--(heap.north-|heap.east) (stack.south-|stack.west)--(stack.south-|stack.east);
  \draw[thick] (static.north-|static.west)--(static.north-|static.east);

  % Arrows
  \draw[draw=CSUcyan, line width=3pt, ->] ([yshift=2pt] stack.south) -- ([yshift=-0.4cm]stack.south);
  \draw[draw=CSUcyan, line width=3pt, ->] ([yshift=-2pt] heap.north) -- ([yshift=0.4cm]heap.north);

  % Numbers
  \node[address, anchor=east] (0) at (static.south west) {$0$};
  \node[address, anchor=south east] (1) at (0.north east) {$1$};
  \node[address, anchor=south] (2) at (1.north) {$\vdots$};

  \node[address, anchor=north east] (h_prev) at (heap.north west) {$h-1$};
  \node[address, anchor=south east] (h) at (heap.north west) {$h$};
  \node[address, anchor=north east] (a) at (stack.south west) {$a$};
  \node[address, anchor=south east] (a_prev) at (stack.south west) {$a - 1$};

  \node[address, anchor=east] (n) at (stack.north west) {$n$};

\end{tikzpicture}
\end{columns}

\end{frame}

\begin{frame}{Static Data Memory}
\begin{columns}
\column{0.67\textwidth}
  \begin{itemize}
  \item<+->
    The simplest memory to manage.
  \item<+->
    Consists of anything that can be completely determined at compile
    time. For example:
    \begin{itemize}
      \item
        global variables,
      \item
        static variables, and
      \item
        function and class definitions (instructions).
    \end{itemize}
  \item<+->
    Characteristics:

    \begin{itemize}
    \item
      Storage requirements are known before execution.
    \item
      The size of the static storage area is constant throughout execution.
    \end{itemize}
  \end{itemize}
\column{0.33\textwidth}
  % % A table representing the layout of runtime memory.
% % Requires \usepackage{arydshln} % for dashed lines in the table
% % Requires \usepackage{hhline} for horizontal lines that show up with \columncolor
% \begin{tabular}{r |>{\columncolor{CSUslate}}c|}
% % \begin{tabular}{r |c|}
%   %\cline{2-2}
%   \hhline{~|-|}
%   $0$ & \textbf{Static}\\
%   $1$ & \\
%       & \\\hhline{~-}
%       & \textbf{Stack}\\
%   $a - 1$ & \\\hhline{~-}
%   \multicolumn{1}{r:}{$a$} &
%     \multicolumn{1}{>{\columncolor{CSUslate}}c:}{$\Downarrow$}\\
%     \multicolumn{1}{r:}{}  & \multicolumn{1}{>{\columncolor{CSUslate}}r:}{} \\\hhline{~-}
%   $h$ & \textbf{Heap}\\
%       &  \\
%   $n$ & \\\hhline{~-}
% \end{tabular}

\begin{tikzpicture}[memory_type/.style={minimum width=2.5cm, minimum height=1.25cm, fill=CSUcyan, text=black}, line/.style={dashed, very thick}, address/.style={font=\small, inner ysep=1pt}]
  % % Define colors for sections
  % \definecolor{staticcolor}{RGB}{173, 216, 230} % Light blue
  % \definecolor{stackcolor}{RGB}{255, 182, 193}  % Light pink
  % \definecolor{heapcolor}{RGB}{144, 238, 144}   % Light green

  % Static Memory Section
  \node[memory_type, fill=CSUcyan!50!white] (static) at (0, 0) {Static};

  % Heap Node (above Static)
  \node[memory_type, anchor=south] (heap) at (static.north) {Heap};

  % Free Node (above Heap)
  \node[memory_type, fill=none, text=white, anchor=south, align=center, font={\small\itshape}, minimum height=1.5cm] (Free) at (heap.north) {unallocated};

  % Stack Node (above Heap)
  \node[memory_type, anchor=south, align=center] (stack) at (Free.north) {Stack};

  \node[fit=(static) (stack), draw, inner sep=0pt, very thick] (Memory) {};

  % % Dashed lines where stack and heap grow.
  % \draw[line] (heap.north-|heap.west)--(heap.north-|heap.east) (stack.south-|stack.west)--(stack.south-|stack.east);
  \draw[thick] (static.north-|static.west)--(static.north-|static.east);

  % Arrows
  \draw[draw=CSUcyan, line width=3pt, ->] ([yshift=2pt] stack.south) -- ([yshift=-0.4cm]stack.south);
  \draw[draw=CSUcyan, line width=3pt, ->] ([yshift=-2pt] heap.north) -- ([yshift=0.4cm]heap.north);

  % Numbers
  \node[address, anchor=east] (0) at (static.south west) {$0$};
  \node[address, anchor=south east] (1) at (0.north east) {$1$};
  \node[address, anchor=south] (2) at (1.north) {$\vdots$};

  \node[address, anchor=north east] (h_prev) at (heap.north west) {$h-1$};
  \node[address, anchor=south east] (h) at (heap.north west) {$h$};
  \node[address, anchor=north east] (a) at (stack.south west) {$a$};
  \node[address, anchor=south east] (a_prev) at (stack.south west) {$a - 1$};

  \node[address, anchor=east] (n) at (stack.north west) {$n$};

\end{tikzpicture}
\end{columns}
\end{frame}

\section{The Stack}


\begin{frame}{The Call Stack}
\begin{columns}
  \column{0.67\textwidth}
  \begin{itemize}[<+->]
  \item
    The \emph{Call Stack} (i.e., the \textit{Runtime Stack} or \textit{Execution Stack}) is a contiguous memory region that grows and shrinks.
  \item
    Supports function calls.
  \item
    The stack \textbf{grows} when a function is called (activated); its \emph{stack frame} 
    (or \emph{activation record}) is pushed on top.
    % It \textbf{grows} (storage is allocated) when the \emph{activation
    % record} (or \emph{stack frame}) is pushed on the stack when a
    % function is called (activated).
  \item
    It \textbf{shrinks} when the function terminates; storage is
    deallocated.
  \end{itemize}

\column{0.33\textwidth}
  \vspace{1em}
  % % A table representing the layout of runtime memory.
% % Requires \usepackage{arydshln} % for dashed lines in the table
% % Requires \usepackage{hhline} for horizontal lines that show up with \columncolor
% \begin{tabular}{r |>{\columncolor{CSUslate}}c|}
% % \begin{tabular}{r |c|}
%   %\cline{2-2}
%   \hhline{~|-|}
%   $0$ & \textbf{Static}\\
%   $1$ & \\
%       & \\\hhline{~-}
%       & \textbf{Stack}\\
%   $a - 1$ & \\\hhline{~-}
%   \multicolumn{1}{r:}{$a$} &
%     \multicolumn{1}{>{\columncolor{CSUslate}}c:}{$\Downarrow$}\\
%     \multicolumn{1}{r:}{}  & \multicolumn{1}{>{\columncolor{CSUslate}}r:}{} \\\hhline{~-}
%   $h$ & \textbf{Heap}\\
%       &  \\
%   $n$ & \\\hhline{~-}
% \end{tabular}

\begin{tikzpicture}[memory_type/.style={minimum width=2.5cm, minimum height=1.25cm, fill=CSUcyan, text=black}, line/.style={dashed, very thick}, address/.style={font=\small, inner ysep=1pt}]
  % % Define colors for sections
  % \definecolor{staticcolor}{RGB}{173, 216, 230} % Light blue
  % \definecolor{stackcolor}{RGB}{255, 182, 193}  % Light pink
  % \definecolor{heapcolor}{RGB}{144, 238, 144}   % Light green

  % Static Memory Section
  \node[memory_type, fill=CSUcyan!50!white] (static) at (0, 0) {Static};

  % Heap Node (above Static)
  \node[memory_type, anchor=south] (heap) at (static.north) {Heap};

  % Free Node (above Heap)
  \node[memory_type, fill=none, text=white, anchor=south, align=center, font={\small\itshape}, minimum height=1.5cm] (Free) at (heap.north) {unallocated};

  % Stack Node (above Heap)
  \node[memory_type, anchor=south, align=center] (stack) at (Free.north) {Stack};

  \node[fit=(static) (stack), draw, inner sep=0pt, very thick] (Memory) {};

  % % Dashed lines where stack and heap grow.
  % \draw[line] (heap.north-|heap.west)--(heap.north-|heap.east) (stack.south-|stack.west)--(stack.south-|stack.east);
  \draw[thick] (static.north-|static.west)--(static.north-|static.east);

  % Arrows
  \draw[draw=CSUcyan, line width=3pt, ->] ([yshift=2pt] stack.south) -- ([yshift=-0.4cm]stack.south);
  \draw[draw=CSUcyan, line width=3pt, ->] ([yshift=-2pt] heap.north) -- ([yshift=0.4cm]heap.north);

  % Numbers
  \node[address, anchor=east] (0) at (static.south west) {$0$};
  \node[address, anchor=south east] (1) at (0.north east) {$1$};
  \node[address, anchor=south] (2) at (1.north) {$\vdots$};

  \node[address, anchor=north east] (h_prev) at (heap.north west) {$h-1$};
  \node[address, anchor=south east] (h) at (heap.north west) {$h$};
  \node[address, anchor=north east] (a) at (stack.south west) {$a$};
  \node[address, anchor=south east] (a_prev) at (stack.south west) {$a - 1$};

  \node[address, anchor=east] (n) at (stack.north west) {$n$};

\end{tikzpicture}
\end{columns}
\end{frame}

\begin{frame}{Stack Frames}
\begin{columns}
\column{0.67\textwidth}
  \begin{itemize}
  \item
    For each function call, a \emph{stack frame} stores local variables, parameters, and
    return linkage.
  \item
    The size and structure of a stack frame are \emph{known at compile time}, but its 
    contents and time of allocation are unknown until runtime.
  \item
    How is variable lifetime affected by stack management techniques?
  \end{itemize}

\column{0.33\textwidth}
  \vspace{1em}
  % % A table representing the layout of runtime memory.
% % Requires \usepackage{arydshln} % for dashed lines in the table
% % Requires \usepackage{hhline} for horizontal lines that show up with \columncolor
% \begin{tabular}{r |>{\columncolor{CSUslate}}c|}
% % \begin{tabular}{r |c|}
%   %\cline{2-2}
%   \hhline{~|-|}
%   $0$ & \textbf{Static}\\
%   $1$ & \\
%       & \\\hhline{~-}
%       & \textbf{Stack}\\
%   $a - 1$ & \\\hhline{~-}
%   \multicolumn{1}{r:}{$a$} &
%     \multicolumn{1}{>{\columncolor{CSUslate}}c:}{$\Downarrow$}\\
%     \multicolumn{1}{r:}{}  & \multicolumn{1}{>{\columncolor{CSUslate}}r:}{} \\\hhline{~-}
%   $h$ & \textbf{Heap}\\
%       &  \\
%   $n$ & \\\hhline{~-}
% \end{tabular}

\begin{tikzpicture}[memory_type/.style={minimum width=2.5cm, minimum height=1.25cm, fill=CSUcyan, text=black}, line/.style={dashed, very thick}, address/.style={font=\small, inner ysep=1pt}]
  % % Define colors for sections
  % \definecolor{staticcolor}{RGB}{173, 216, 230} % Light blue
  % \definecolor{stackcolor}{RGB}{255, 182, 193}  % Light pink
  % \definecolor{heapcolor}{RGB}{144, 238, 144}   % Light green

  % Static Memory Section
  \node[memory_type, fill=CSUcyan!50!white] (static) at (0, 0) {Static};

  % Heap Node (above Static)
  \node[memory_type, anchor=south] (heap) at (static.north) {Heap};

  % Free Node (above Heap)
  \node[memory_type, fill=none, text=white, anchor=south, align=center, font={\small\itshape}, minimum height=1.5cm] (Free) at (heap.north) {unallocated};

  % Stack Node (above Heap)
  \node[memory_type, anchor=south, align=center] (stack) at (Free.north) {Stack};

  \node[fit=(static) (stack), draw, inner sep=0pt, very thick] (Memory) {};

  % % Dashed lines where stack and heap grow.
  % \draw[line] (heap.north-|heap.west)--(heap.north-|heap.east) (stack.south-|stack.west)--(stack.south-|stack.east);
  \draw[thick] (static.north-|static.west)--(static.north-|static.east);

  % Arrows
  \draw[draw=CSUcyan, line width=3pt, ->] ([yshift=2pt] stack.south) -- ([yshift=-0.4cm]stack.south);
  \draw[draw=CSUcyan, line width=3pt, ->] ([yshift=-2pt] heap.north) -- ([yshift=0.4cm]heap.north);

  % Numbers
  \node[address, anchor=east] (0) at (static.south west) {$0$};
  \node[address, anchor=south east] (1) at (0.north east) {$1$};
  \node[address, anchor=south] (2) at (1.north) {$\vdots$};

  \node[address, anchor=north east] (h_prev) at (heap.north west) {$h-1$};
  \node[address, anchor=south east] (h) at (heap.north west) {$h$};
  \node[address, anchor=north east] (a) at (stack.south west) {$a$};
  \node[address, anchor=south east] (a_prev) at (stack.south west) {$a - 1$};

  \node[address, anchor=east] (n) at (stack.north west) {$n$};

\end{tikzpicture}
\end{columns}
\end{frame}

\begin{frame}[fragile]{Stack Frames}

\href{https://pythontutor.com/render.html#code=int%20add%28int%20a,%20int%20b%29%0A%7B%0A%20%20return%20a%20%2B%20b%3B%0A%7D%0A%0Aint%20doubleAdd%28int%20a,%20int%20b%29%0A%7B%0A%20%20return%20add%28a,%20b%29%0A%20%20%20%20%2B%20add%28b,%20a%29%3B%0A%7D%0A%0Aint%20main%28%29%0A%7B%0A%20%20int%20total%20%3D%200%3B%0A%20%20total%20%3D%20add%281,%202%29%3B%0A%20%20total%20%3D%20doubleAdd%283,%204%29%3B%0A%0A%20%20return%200%3B%0A%7D&cumulative=false&curInstr=0&heapPrimitives=nevernest&mode=display&origin=opt-frontend.js&py=cpp_g%2B%2B9.3.0&rawInputLstJSON=%5B%5D&textReferences=false}{Click here to step through this example code.}

\begin{columns}

\column{0.67\textwidth}

\begin{lstlisting}[basicstyle=\ttfamily\bfseries\scriptsize,frame=single,numbers=left,
  numbersep=5pt,
  numberstyle=\tiny\color{gray},
  xleftmargin=20pt]
int add(int a, int b)
{
  return a + b;  // <- 3
}

int doubleAdd(int a, int b)
{
  return add(a, b)
    + add(b, a);  // <- 2
}

int main()
{
  int total = 0;
  total = add(1, 2);
  total = doubleAdd(3, 4); // <- 1

  return 0;
}
\end{lstlisting}

\column{0.33\textwidth}
\begin{tabular}{r |>{\ttfamily\columncolor{CSUslate}}l|}
  \multicolumn{1}{r}{} & \multicolumn{1}{c}{\textbf{Call Stack}}\\
    \hhline{~|-|}%\cline{2-2}
      & \texttt{main()}\\
      & total $\rightarrow$ 0\\
    \hhline{~|-|}%\cline{2-2}
      & \texttt{doubleAdd()}\\
      & int a $\rightarrow$ 3\\
      & int b $\rightarrow$ 4\\
    \hhline{~|-|}%\cline{2-2}
      & \texttt{add()}\\
      & int a $\rightarrow$ 4 \\
$sp$ $\rightarrow$ & int b $\rightarrow$ 3 \\
  \cline{2-2}
\end{tabular}
\end{columns}
\end{frame}

\begin{frame}{Stack Overflow}
\begin{columns}
\column{0.67\textwidth}

  \begin{itemize}
  \item
    The call stack and heap grow towards each other as required by program events.
  \item
    The following relation must hold:\\
    $0 \le h \le a \le n$
  \item
    In other words, if the stack top bumps into the heap, or if the beginning of the heap is greater than the end, there are problems!
  \end{itemize}

\column{0.33\textwidth}
  % % A table representing the layout of runtime memory.
% % Requires \usepackage{arydshln} % for dashed lines in the table
% % Requires \usepackage{hhline} for horizontal lines that show up with \columncolor
% \begin{tabular}{r |>{\columncolor{CSUslate}}c|}
% % \begin{tabular}{r |c|}
%   %\cline{2-2}
%   \hhline{~|-|}
%   $0$ & \textbf{Static}\\
%   $1$ & \\
%       & \\\hhline{~-}
%       & \textbf{Stack}\\
%   $a - 1$ & \\\hhline{~-}
%   \multicolumn{1}{r:}{$a$} &
%     \multicolumn{1}{>{\columncolor{CSUslate}}c:}{$\Downarrow$}\\
%     \multicolumn{1}{r:}{}  & \multicolumn{1}{>{\columncolor{CSUslate}}r:}{} \\\hhline{~-}
%   $h$ & \textbf{Heap}\\
%       &  \\
%   $n$ & \\\hhline{~-}
% \end{tabular}

\begin{tikzpicture}[memory_type/.style={minimum width=2.5cm, minimum height=1.25cm, fill=CSUcyan, text=black}, line/.style={dashed, very thick}, address/.style={font=\small, inner ysep=1pt}]
  % % Define colors for sections
  % \definecolor{staticcolor}{RGB}{173, 216, 230} % Light blue
  % \definecolor{stackcolor}{RGB}{255, 182, 193}  % Light pink
  % \definecolor{heapcolor}{RGB}{144, 238, 144}   % Light green

  % Static Memory Section
  \node[memory_type, fill=CSUcyan!50!white] (static) at (0, 0) {Static};

  % Heap Node (above Static)
  \node[memory_type, anchor=south] (heap) at (static.north) {Heap};

  % Free Node (above Heap)
  \node[memory_type, fill=none, text=white, anchor=south, align=center, font={\small\itshape}, minimum height=1.5cm] (Free) at (heap.north) {unallocated};

  % Stack Node (above Heap)
  \node[memory_type, anchor=south, align=center] (stack) at (Free.north) {Stack};

  \node[fit=(static) (stack), draw, inner sep=0pt, very thick] (Memory) {};

  % % Dashed lines where stack and heap grow.
  % \draw[line] (heap.north-|heap.west)--(heap.north-|heap.east) (stack.south-|stack.west)--(stack.south-|stack.east);
  \draw[thick] (static.north-|static.west)--(static.north-|static.east);

  % Arrows
  \draw[draw=CSUcyan, line width=3pt, ->] ([yshift=2pt] stack.south) -- ([yshift=-0.4cm]stack.south);
  \draw[draw=CSUcyan, line width=3pt, ->] ([yshift=-2pt] heap.north) -- ([yshift=0.4cm]heap.north);

  % Numbers
  \node[address, anchor=east] (0) at (static.south west) {$0$};
  \node[address, anchor=south east] (1) at (0.north east) {$1$};
  \node[address, anchor=south] (2) at (1.north) {$\vdots$};

  \node[address, anchor=north east] (h_prev) at (heap.north west) {$h-1$};
  \node[address, anchor=south east] (h) at (heap.north west) {$h$};
  \node[address, anchor=north east] (a) at (stack.south west) {$a$};
  \node[address, anchor=south east] (a_prev) at (stack.south west) {$a - 1$};

  \node[address, anchor=east] (n) at (stack.north west) {$n$};

\end{tikzpicture}
\end{columns}

\end{frame}

\section{The Heap}

\begin{frame}{Heap Memory}
\begin{columns}[b]
\column{0.67\textwidth}
  \begin{itemize}[<+->]
  \item
    Heap objects are \textit{dynamically} allocated/deallocated at runtime
    (not associated with function call/return).
  \item
    Dynamic variables are created at runtime instead of at compile-time.
  \item
    The kind of data found on the heap is language dependant.

    \begin{itemize}
    \item
      Typically holds strings, dynamic arrays, objects, and linked structures
    \item
      Java and C/C++ have different policies.
    \end{itemize}
  \end{itemize}
\column{0.3\textwidth}
  \vfill{}
  % % A table representing the layout of runtime memory.
% % Requires \usepackage{arydshln} % for dashed lines in the table
% % Requires \usepackage{hhline} for horizontal lines that show up with \columncolor
% \begin{tabular}{r |>{\columncolor{CSUslate}}c|}
% % \begin{tabular}{r |c|}
%   %\cline{2-2}
%   \hhline{~|-|}
%   $0$ & \textbf{Static}\\
%   $1$ & \\
%       & \\\hhline{~-}
%       & \textbf{Stack}\\
%   $a - 1$ & \\\hhline{~-}
%   \multicolumn{1}{r:}{$a$} &
%     \multicolumn{1}{>{\columncolor{CSUslate}}c:}{$\Downarrow$}\\
%     \multicolumn{1}{r:}{}  & \multicolumn{1}{>{\columncolor{CSUslate}}r:}{} \\\hhline{~-}
%   $h$ & \textbf{Heap}\\
%       &  \\
%   $n$ & \\\hhline{~-}
% \end{tabular}

\begin{tikzpicture}[memory_type/.style={minimum width=2.5cm, minimum height=1.25cm, fill=CSUcyan, text=black}, line/.style={dashed, very thick}, address/.style={font=\small, inner ysep=1pt}]
  % % Define colors for sections
  % \definecolor{staticcolor}{RGB}{173, 216, 230} % Light blue
  % \definecolor{stackcolor}{RGB}{255, 182, 193}  % Light pink
  % \definecolor{heapcolor}{RGB}{144, 238, 144}   % Light green

  % Static Memory Section
  \node[memory_type, fill=CSUcyan!50!white] (static) at (0, 0) {Static};

  % Heap Node (above Static)
  \node[memory_type, anchor=south] (heap) at (static.north) {Heap};

  % Free Node (above Heap)
  \node[memory_type, fill=none, text=white, anchor=south, align=center, font={\small\itshape}, minimum height=1.5cm] (Free) at (heap.north) {unallocated};

  % Stack Node (above Heap)
  \node[memory_type, anchor=south, align=center] (stack) at (Free.north) {Stack};

  \node[fit=(static) (stack), draw, inner sep=0pt, very thick] (Memory) {};

  % % Dashed lines where stack and heap grow.
  % \draw[line] (heap.north-|heap.west)--(heap.north-|heap.east) (stack.south-|stack.west)--(stack.south-|stack.east);
  \draw[thick] (static.north-|static.west)--(static.north-|static.east);

  % Arrows
  \draw[draw=CSUcyan, line width=3pt, ->] ([yshift=2pt] stack.south) -- ([yshift=-0.4cm]stack.south);
  \draw[draw=CSUcyan, line width=3pt, ->] ([yshift=-2pt] heap.north) -- ([yshift=0.4cm]heap.north);

  % Numbers
  \node[address, anchor=east] (0) at (static.south west) {$0$};
  \node[address, anchor=south east] (1) at (0.north east) {$1$};
  \node[address, anchor=south] (2) at (1.north) {$\vdots$};

  \node[address, anchor=north east] (h_prev) at (heap.north west) {$h-1$};
  \node[address, anchor=south east] (h) at (heap.north west) {$h$};
  \node[address, anchor=north east] (a) at (stack.south west) {$a$};
  \node[address, anchor=south east] (a_prev) at (stack.south west) {$a - 1$};

  \node[address, anchor=east] (n) at (stack.north west) {$n$};

\end{tikzpicture}
\end{columns}
\end{frame}


\begin{frame}[fragile]{Heap Memory Example}

\href{https://pythontutor.com/render.html#code=%23include%20%3Ciostream%3E%0A%0Aint%20main%28%29%0A%7B%0A%20%20const%20int%20SIZE%20%3D%203%3B%0A%20%20int%20stackArray%5BSIZE%5D%3B%20//%20Declared%20on%20the%20stack%0A%20%20int*%20heapArray%3B%20//%20Pointer%20to%20memory%20location%0A%20%20heapArray%20%3D%20new%20int%5BSIZE%5D%3B%20//%20Declare%20array%20on%20heap%0A%0A%20%20std%3A%3Acout%20%3C%3C%20%22stackArray%20%20%3D%20%22%20%3C%3C%20stackArray%20%3C%3C%20std%3A%3Aendl%3B%0A%20%20std%3A%3Acout%20%3C%3C%20%22heapArray%20%20%20%3D%20%22%20%3C%3C%20heapArray%20%3C%3C%20std%3A%3Aendl%3B%0A%20%20%0A%20%20stackArray%5B2%5D%20%3D%2020%3B%0A%20%20heapArray%5B2%5D%20%3D%2020%3B%0A%20%20%0A%20%20//%20Free%20up%20the%20memory%20from%20the%20heap%20array%0A%20%20delete%5B%5D%20heapArray%3B%0A%20%20heapArray%20%3D%20nullptr%3B%0A%0A%20%20return%200%3B%0A%7D&cumulative=false&curInstr=0&heapPrimitives=nevernest&mode=display&origin=opt-frontend.js&py=cpp_g%2B%2B9.3.0&rawInputLstJSON=%5B%5D&textReferences=false}{Click here to step through this example code.}%
\begin{columns}%
\column{0.9\textwidth}%
\begin{lstlisting}[basicstyle=\ttfamily\bfseries\scriptsize,frame=single,
  numbers=left, numbersep=5pt,numberstyle=\tiny\color{gray},
  xleftmargin=20pt]
#include <iostream>

int main()
{
  const int SIZE = 3;
  int stackArray[SIZE];      // Declared on the stack
  int *heapArray;            // Declare pointer to memory location
  heapArray = new int[SIZE]; // Declare array on heap

  std::cout << "stackArray  = " << stackArray << std::endl;
  std::cout << "heapArray   = " << heapArray << std::endl;

  stackArray[2] = 20;
  heapArray[2] = 20;

  delete[] heapArray; // Free up the memory from the heap array
  heapArray = nullptr; // Set pointer to point to address 0.

  return 0;
}
\end{lstlisting}

\column{0.1\textwidth}

\end{columns}
\end{frame}

\begin{frame}[fragile]{Heap Memory}
\begin{columns}
\column{0.67\textwidth}
  \begin{itemize}
  \item
    The \lstinline{new} operator allocates heap storage.
  \item
    The \lstinline{delete} or \lstinline{delete[]} operators deallocate heap storage for reuse.
  \item
    Space is allocated in variable-sized blocks, so deallocation may leave
    ``holes'' in the heap (fragmentation).

    \begin{itemize}
    \item
      Compared to the deallocation of stack storage
    \end{itemize}
  \end{itemize}

\column{0.3\textwidth}
  % % A table representing the layout of runtime memory.
% % Requires \usepackage{arydshln} % for dashed lines in the table
% % Requires \usepackage{hhline} for horizontal lines that show up with \columncolor
% \begin{tabular}{r |>{\columncolor{CSUslate}}c|}
% % \begin{tabular}{r |c|}
%   %\cline{2-2}
%   \hhline{~|-|}
%   $0$ & \textbf{Static}\\
%   $1$ & \\
%       & \\\hhline{~-}
%       & \textbf{Stack}\\
%   $a - 1$ & \\\hhline{~-}
%   \multicolumn{1}{r:}{$a$} &
%     \multicolumn{1}{>{\columncolor{CSUslate}}c:}{$\Downarrow$}\\
%     \multicolumn{1}{r:}{}  & \multicolumn{1}{>{\columncolor{CSUslate}}r:}{} \\\hhline{~-}
%   $h$ & \textbf{Heap}\\
%       &  \\
%   $n$ & \\\hhline{~-}
% \end{tabular}

\begin{tikzpicture}[memory_type/.style={minimum width=2.5cm, minimum height=1.25cm, fill=CSUcyan, text=black}, line/.style={dashed, very thick}, address/.style={font=\small, inner ysep=1pt}]
  % % Define colors for sections
  % \definecolor{staticcolor}{RGB}{173, 216, 230} % Light blue
  % \definecolor{stackcolor}{RGB}{255, 182, 193}  % Light pink
  % \definecolor{heapcolor}{RGB}{144, 238, 144}   % Light green

  % Static Memory Section
  \node[memory_type, fill=CSUcyan!50!white] (static) at (0, 0) {Static};

  % Heap Node (above Static)
  \node[memory_type, anchor=south] (heap) at (static.north) {Heap};

  % Free Node (above Heap)
  \node[memory_type, fill=none, text=white, anchor=south, align=center, font={\small\itshape}, minimum height=1.5cm] (Free) at (heap.north) {unallocated};

  % Stack Node (above Heap)
  \node[memory_type, anchor=south, align=center] (stack) at (Free.north) {Stack};

  \node[fit=(static) (stack), draw, inner sep=0pt, very thick] (Memory) {};

  % % Dashed lines where stack and heap grow.
  % \draw[line] (heap.north-|heap.west)--(heap.north-|heap.east) (stack.south-|stack.west)--(stack.south-|stack.east);
  \draw[thick] (static.north-|static.west)--(static.north-|static.east);

  % Arrows
  \draw[draw=CSUcyan, line width=3pt, ->] ([yshift=2pt] stack.south) -- ([yshift=-0.4cm]stack.south);
  \draw[draw=CSUcyan, line width=3pt, ->] ([yshift=-2pt] heap.north) -- ([yshift=0.4cm]heap.north);

  % Numbers
  \node[address, anchor=east] (0) at (static.south west) {$0$};
  \node[address, anchor=south east] (1) at (0.north east) {$1$};
  \node[address, anchor=south] (2) at (1.north) {$\vdots$};

  \node[address, anchor=north east] (h_prev) at (heap.north west) {$h-1$};
  \node[address, anchor=south east] (h) at (heap.north west) {$h$};
  \node[address, anchor=north east] (a) at (stack.south west) {$a$};
  \node[address, anchor=south east] (a_prev) at (stack.south west) {$a - 1$};

  \node[address, anchor=east] (n) at (stack.north west) {$n$};

\end{tikzpicture}
\end{columns}
\end{frame}

\begin{frame}[fragile]{Heap Management}
\begin{columns}
\column{0.67\textwidth}
  \begin{itemize}[<+->]
  \item
    Some languages (e.g., C, C++) leave heap storage deallocation to the
    programmer.
  \item
    Others (e.g., Java, Perl, Python, list-processing languages) employ
    \textit{garbage collection} to reclaim unused heap space.
  \end{itemize}
\column{0.3\textwidth}
  % % A table representing the layout of runtime memory.
% % Requires \usepackage{arydshln} % for dashed lines in the table
% % Requires \usepackage{hhline} for horizontal lines that show up with \columncolor
% \begin{tabular}{r |>{\columncolor{CSUslate}}c|}
% % \begin{tabular}{r |c|}
%   %\cline{2-2}
%   \hhline{~|-|}
%   $0$ & \textbf{Static}\\
%   $1$ & \\
%       & \\\hhline{~-}
%       & \textbf{Stack}\\
%   $a - 1$ & \\\hhline{~-}
%   \multicolumn{1}{r:}{$a$} &
%     \multicolumn{1}{>{\columncolor{CSUslate}}c:}{$\Downarrow$}\\
%     \multicolumn{1}{r:}{}  & \multicolumn{1}{>{\columncolor{CSUslate}}r:}{} \\\hhline{~-}
%   $h$ & \textbf{Heap}\\
%       &  \\
%   $n$ & \\\hhline{~-}
% \end{tabular}

\begin{tikzpicture}[memory_type/.style={minimum width=2.5cm, minimum height=1.25cm, fill=CSUcyan, text=black}, line/.style={dashed, very thick}, address/.style={font=\small, inner ysep=1pt}]
  % % Define colors for sections
  % \definecolor{staticcolor}{RGB}{173, 216, 230} % Light blue
  % \definecolor{stackcolor}{RGB}{255, 182, 193}  % Light pink
  % \definecolor{heapcolor}{RGB}{144, 238, 144}   % Light green

  % Static Memory Section
  \node[memory_type, fill=CSUcyan!50!white] (static) at (0, 0) {Static};

  % Heap Node (above Static)
  \node[memory_type, anchor=south] (heap) at (static.north) {Heap};

  % Free Node (above Heap)
  \node[memory_type, fill=none, text=white, anchor=south, align=center, font={\small\itshape}, minimum height=1.5cm] (Free) at (heap.north) {unallocated};

  % Stack Node (above Heap)
  \node[memory_type, anchor=south, align=center] (stack) at (Free.north) {Stack};

  \node[fit=(static) (stack), draw, inner sep=0pt, very thick] (Memory) {};

  % % Dashed lines where stack and heap grow.
  % \draw[line] (heap.north-|heap.west)--(heap.north-|heap.east) (stack.south-|stack.west)--(stack.south-|stack.east);
  \draw[thick] (static.north-|static.west)--(static.north-|static.east);

  % Arrows
  \draw[draw=CSUcyan, line width=3pt, ->] ([yshift=2pt] stack.south) -- ([yshift=-0.4cm]stack.south);
  \draw[draw=CSUcyan, line width=3pt, ->] ([yshift=-2pt] heap.north) -- ([yshift=0.4cm]heap.north);

  % Numbers
  \node[address, anchor=east] (0) at (static.south west) {$0$};
  \node[address, anchor=south east] (1) at (0.north east) {$1$};
  \node[address, anchor=south] (2) at (1.north) {$\vdots$};

  \node[address, anchor=north east] (h_prev) at (heap.north west) {$h-1$};
  \node[address, anchor=south east] (h) at (heap.north west) {$h$};
  \node[address, anchor=north east] (a) at (stack.south west) {$a$};
  \node[address, anchor=south east] (a_prev) at (stack.south west) {$a - 1$};

  \node[address, anchor=east] (n) at (stack.north west) {$n$};

\end{tikzpicture}
\end{columns}
\end{frame}

\begin{frame}{Heap Management Functions}
\begin{columns}
\column{0.67\textwidth}
  \begin{itemize}[<+->]
  \item
    \lstinline{new} returns the start address of an unused block from the heap and changes its state from \textit{unused} to \textit{reserved} (and undefined).
  \item
    Suppose a Point class has three 4-byte data members: $x$, $y$, $z$.\\
  \item
    \lstinline[language=Java]{Point firstCoord = new Point();}\\
    requires at least allocated $3 \times 4$ bytes.
  \end{itemize}
\column{0.3\textwidth}
  % % A table representing the layout of runtime memory.
% % Requires \usepackage{arydshln} % for dashed lines in the table
% % Requires \usepackage{hhline} for horizontal lines that show up with \columncolor
% \begin{tabular}{r |>{\columncolor{CSUslate}}c|}
% % \begin{tabular}{r |c|}
%   %\cline{2-2}
%   \hhline{~|-|}
%   $0$ & \textbf{Static}\\
%   $1$ & \\
%       & \\\hhline{~-}
%       & \textbf{Stack}\\
%   $a - 1$ & \\\hhline{~-}
%   \multicolumn{1}{r:}{$a$} &
%     \multicolumn{1}{>{\columncolor{CSUslate}}c:}{$\Downarrow$}\\
%     \multicolumn{1}{r:}{}  & \multicolumn{1}{>{\columncolor{CSUslate}}r:}{} \\\hhline{~-}
%   $h$ & \textbf{Heap}\\
%       &  \\
%   $n$ & \\\hhline{~-}
% \end{tabular}

\begin{tikzpicture}[memory_type/.style={minimum width=2.5cm, minimum height=1.25cm, fill=CSUcyan, text=black}, line/.style={dashed, very thick}, address/.style={font=\small, inner ysep=1pt}]
  % % Define colors for sections
  % \definecolor{staticcolor}{RGB}{173, 216, 230} % Light blue
  % \definecolor{stackcolor}{RGB}{255, 182, 193}  % Light pink
  % \definecolor{heapcolor}{RGB}{144, 238, 144}   % Light green

  % Static Memory Section
  \node[memory_type, fill=CSUcyan!50!white] (static) at (0, 0) {Static};

  % Heap Node (above Static)
  \node[memory_type, anchor=south] (heap) at (static.north) {Heap};

  % Free Node (above Heap)
  \node[memory_type, fill=none, text=white, anchor=south, align=center, font={\small\itshape}, minimum height=1.5cm] (Free) at (heap.north) {unallocated};

  % Stack Node (above Heap)
  \node[memory_type, anchor=south, align=center] (stack) at (Free.north) {Stack};

  \node[fit=(static) (stack), draw, inner sep=0pt, very thick] (Memory) {};

  % % Dashed lines where stack and heap grow.
  % \draw[line] (heap.north-|heap.west)--(heap.north-|heap.east) (stack.south-|stack.west)--(stack.south-|stack.east);
  \draw[thick] (static.north-|static.west)--(static.north-|static.east);

  % Arrows
  \draw[draw=CSUcyan, line width=3pt, ->] ([yshift=2pt] stack.south) -- ([yshift=-0.4cm]stack.south);
  \draw[draw=CSUcyan, line width=3pt, ->] ([yshift=-2pt] heap.north) -- ([yshift=0.4cm]heap.north);

  % Numbers
  \node[address, anchor=east] (0) at (static.south west) {$0$};
  \node[address, anchor=south east] (1) at (0.north east) {$1$};
  \node[address, anchor=south] (2) at (1.north) {$\vdots$};

  \node[address, anchor=north east] (h_prev) at (heap.north west) {$h-1$};
  \node[address, anchor=south east] (h) at (heap.north west) {$h$};
  \node[address, anchor=north east] (a) at (stack.south west) {$a$};
  \node[address, anchor=south east] (a_prev) at (stack.south west) {$a - 1$};

  \node[address, anchor=east] (n) at (stack.north west) {$n$};

\end{tikzpicture}
\end{columns}
\end{frame}

\begin{frame}{Heap Overflow}
\begin{columns}
\column{0.67\textwidth}
  \begin{itemize}[<+->]
  \item
    A \textbf{heap overflow} occurs when calling \lstinline{new} and the heap
    does not have a big enough block of free contiguous space.
  \item
    So, \lstinline{new} either fails (in the case of heap overflow) or returns a
    pointer to the new block.
  \end{itemize}
\column{0.3\textwidth}
  % % A table representing the layout of runtime memory.
% % Requires \usepackage{arydshln} % for dashed lines in the table
% % Requires \usepackage{hhline} for horizontal lines that show up with \columncolor
% \begin{tabular}{r |>{\columncolor{CSUslate}}c|}
% % \begin{tabular}{r |c|}
%   %\cline{2-2}
%   \hhline{~|-|}
%   $0$ & \textbf{Static}\\
%   $1$ & \\
%       & \\\hhline{~-}
%       & \textbf{Stack}\\
%   $a - 1$ & \\\hhline{~-}
%   \multicolumn{1}{r:}{$a$} &
%     \multicolumn{1}{>{\columncolor{CSUslate}}c:}{$\Downarrow$}\\
%     \multicolumn{1}{r:}{}  & \multicolumn{1}{>{\columncolor{CSUslate}}r:}{} \\\hhline{~-}
%   $h$ & \textbf{Heap}\\
%       &  \\
%   $n$ & \\\hhline{~-}
% \end{tabular}

\begin{tikzpicture}[memory_type/.style={minimum width=2.5cm, minimum height=1.25cm, fill=CSUcyan, text=black}, line/.style={dashed, very thick}, address/.style={font=\small, inner ysep=1pt}]
  % % Define colors for sections
  % \definecolor{staticcolor}{RGB}{173, 216, 230} % Light blue
  % \definecolor{stackcolor}{RGB}{255, 182, 193}  % Light pink
  % \definecolor{heapcolor}{RGB}{144, 238, 144}   % Light green

  % Static Memory Section
  \node[memory_type, fill=CSUcyan!50!white] (static) at (0, 0) {Static};

  % Heap Node (above Static)
  \node[memory_type, anchor=south] (heap) at (static.north) {Heap};

  % Free Node (above Heap)
  \node[memory_type, fill=none, text=white, anchor=south, align=center, font={\small\itshape}, minimum height=1.5cm] (Free) at (heap.north) {unallocated};

  % Stack Node (above Heap)
  \node[memory_type, anchor=south, align=center] (stack) at (Free.north) {Stack};

  \node[fit=(static) (stack), draw, inner sep=0pt, very thick] (Memory) {};

  % % Dashed lines where stack and heap grow.
  % \draw[line] (heap.north-|heap.west)--(heap.north-|heap.east) (stack.south-|stack.west)--(stack.south-|stack.east);
  \draw[thick] (static.north-|static.west)--(static.north-|static.east);

  % Arrows
  \draw[draw=CSUcyan, line width=3pt, ->] ([yshift=2pt] stack.south) -- ([yshift=-0.4cm]stack.south);
  \draw[draw=CSUcyan, line width=3pt, ->] ([yshift=-2pt] heap.north) -- ([yshift=0.4cm]heap.north);

  % Numbers
  \node[address, anchor=east] (0) at (static.south west) {$0$};
  \node[address, anchor=south east] (1) at (0.north east) {$1$};
  \node[address, anchor=south] (2) at (1.north) {$\vdots$};

  \node[address, anchor=north east] (h_prev) at (heap.north west) {$h-1$};
  \node[address, anchor=south east] (h) at (heap.north west) {$h$};
  \node[address, anchor=north east] (a) at (stack.south west) {$a$};
  \node[address, anchor=south east] (a_prev) at (stack.south west) {$a - 1$};

  \node[address, anchor=east] (n) at (stack.north west) {$n$};

\end{tikzpicture}
\end{columns}
\end{frame}

\begin{frame}{Heap Management Functions}
\begin{columns}
\column{0.67\textwidth}
  \begin{itemize}[<+->]
  \item
    \lstinline{delete} returns a storage block to the heap.
  \item
    The status of the block returns to \textit{unused} and
    is available for allocation by future calls to \lstinline{new}.
  \item
    One cause of heap overflow is a failure on the part of the program to
    return unused storage.
  \end{itemize}
\column{0.3\textwidth}
  % % A table representing the layout of runtime memory.
% % Requires \usepackage{arydshln} % for dashed lines in the table
% % Requires \usepackage{hhline} for horizontal lines that show up with \columncolor
% \begin{tabular}{r |>{\columncolor{CSUslate}}c|}
% % \begin{tabular}{r |c|}
%   %\cline{2-2}
%   \hhline{~|-|}
%   $0$ & \textbf{Static}\\
%   $1$ & \\
%       & \\\hhline{~-}
%       & \textbf{Stack}\\
%   $a - 1$ & \\\hhline{~-}
%   \multicolumn{1}{r:}{$a$} &
%     \multicolumn{1}{>{\columncolor{CSUslate}}c:}{$\Downarrow$}\\
%     \multicolumn{1}{r:}{}  & \multicolumn{1}{>{\columncolor{CSUslate}}r:}{} \\\hhline{~-}
%   $h$ & \textbf{Heap}\\
%       &  \\
%   $n$ & \\\hhline{~-}
% \end{tabular}

\begin{tikzpicture}[memory_type/.style={minimum width=2.5cm, minimum height=1.25cm, fill=CSUcyan, text=black}, line/.style={dashed, very thick}, address/.style={font=\small, inner ysep=1pt}]
  % % Define colors for sections
  % \definecolor{staticcolor}{RGB}{173, 216, 230} % Light blue
  % \definecolor{stackcolor}{RGB}{255, 182, 193}  % Light pink
  % \definecolor{heapcolor}{RGB}{144, 238, 144}   % Light green

  % Static Memory Section
  \node[memory_type, fill=CSUcyan!50!white] (static) at (0, 0) {Static};

  % Heap Node (above Static)
  \node[memory_type, anchor=south] (heap) at (static.north) {Heap};

  % Free Node (above Heap)
  \node[memory_type, fill=none, text=white, anchor=south, align=center, font={\small\itshape}, minimum height=1.5cm] (Free) at (heap.north) {unallocated};

  % Stack Node (above Heap)
  \node[memory_type, anchor=south, align=center] (stack) at (Free.north) {Stack};

  \node[fit=(static) (stack), draw, inner sep=0pt, very thick] (Memory) {};

  % % Dashed lines where stack and heap grow.
  % \draw[line] (heap.north-|heap.west)--(heap.north-|heap.east) (stack.south-|stack.west)--(stack.south-|stack.east);
  \draw[thick] (static.north-|static.west)--(static.north-|static.east);

  % Arrows
  \draw[draw=CSUcyan, line width=3pt, ->] ([yshift=2pt] stack.south) -- ([yshift=-0.4cm]stack.south);
  \draw[draw=CSUcyan, line width=3pt, ->] ([yshift=-2pt] heap.north) -- ([yshift=0.4cm]heap.north);

  % Numbers
  \node[address, anchor=east] (0) at (static.south west) {$0$};
  \node[address, anchor=south east] (1) at (0.north east) {$1$};
  \node[address, anchor=south] (2) at (1.north) {$\vdots$};

  \node[address, anchor=north east] (h_prev) at (heap.north west) {$h-1$};
  \node[address, anchor=south east] (h) at (heap.north west) {$h$};
  \node[address, anchor=north east] (a) at (stack.south west) {$a$};
  \node[address, anchor=south east] (a_prev) at (stack.south west) {$a - 1$};

  \node[address, anchor=east] (n) at (stack.north west) {$n$};

\end{tikzpicture}
\end{columns}
\end{frame}

\section{Pointers}
\begin{frame}{Pointers}

\begin{itemize}[<+->]
\item
  Pointers are addresses (i.e., the value of a pointer variable is an address).
\item
  Memory that is accessed through a pointer is dynamically allocated in
  the heap.
\item
  Java doesn't have explicit pointers; reference types are heap allocated
   (although the reference is on the stack).
\item
  This topic is covered well in the textbook's Chapter 12.
\end{itemize}
\end{frame}

\begin{frame}{Memory Allocation Java vs C++}
\small
\begin{tabularx}{1\textwidth}{| l l X |}
\hline
\rowcolor{CSUgold}\textcolor{black}{\textbf{Java}}  & \textcolor{black}{Stack Example}   & \textcolor{black}{Heap Example} \\ \hline
Primitives & \cellcolor{codeBackground}\lstinline{int a = 3;}   & \cellcolor{codeBackground}\lstinline[language=Java]{int{[}{]} a = new int{[}1{]};} \\ %\cline{2-4}
Classes    & \textit{Does not exist.} & \cellcolor{codeBackground}\lstinline[language=Java]{String str = new String("Hi!");} \\ \hline
\multicolumn{3}{c}{}\\\hline

\rowcolor{CSUgold}\textcolor{black}{\textbf{C++}}  & \textcolor{black}{Stack Example}   & \textcolor{black}{Heap Example} \\ \hline
Primitives & \cellcolor{codeBackground}\lstinline{int a = 3;}   & \cellcolor{codeBackground}\lstinline{int *a = new int;} \\ %\cline{2-4}
Classes    & \cellcolor{codeBackground}\lstinline{string str("Hi!");}  & \cellcolor{codeBackground}\lstinline{string *str = new string("Hi!");}\\ \hline
\end{tabularx}

\end{frame}

\begin{frame}{Examples in Code}
\centering\Large
Let's create and test out some more pointers.
\end{frame}

\begin{frame}{Java Versus C/C++ Arrays}

\begin{itemize}[<+->]
\item
  In Java, arrays are always allocated dynamically from heap memory.
\item
  In many other languages, including C++:\\
  \begin{tabular}{l @{ -- } l}
    Globally defined arrays & static memory.\\
    Local (to a function) arrays & stack storage.\\
    Dynamically allocated arrays & heap storage.\\
  \end{tabular}
\item
  Dynamically allocated arrays also have stack storage --- a
  reference (pointer) to the heap block that holds the array.
\end{itemize}
\end{frame}

\begin{frame}[fragile]{Comparing Pointers}

\begin{lstlisting}
int i = 5, j = 5;
int *ptrJ1 = &j;
int *ptrJ2 = &j;
int *ptrI = &i;
\end{lstlisting}


True/False:
\begin{itemize}
\item
  \lstinline|if (ptrJ1 == ptrJ2) {}  | \alt<2->{True}{?}
\item <3->
  \lstinline|if (ptrJ1 == ptrI) {}   | \alt<4->{False}{?}
\item <5->
  \lstinline|if (&ptrJ1 == &ptrJ2) {}| \alt<6->{False}{?}
\item <7->
  \lstinline|if (*ptrJ1 == *ptrI) {} | \alt<8->{True}{?}
\end{itemize}

\href{https://pythontutor.com/visualize.html#code=%23include%20%3Ciostream%3E%0A%0Aint%20main%28%29%0A%7B%0A%20%20int%20i%20%3D%205,%20j%20%3D%205%3B%0A%20%20int%20*ptrJ1%20%3D%20%26j%3B%0A%20%20int%20*ptrJ2%20%3D%20%26j%3B%0A%20%20int%20*ptrI%20%3D%20%26i%3B%0A%20%20%0A%20%20if%20%28ptrJ1%20%3D%3D%20ptrJ2%29%20%7B%0A%20%20%20%20std%3A%3Acout%20%3C%3C%201%20%3C%3C%20'%20'%3B%0A%20%20%7D%0A%20%20if%20%28ptrJ1%20%3D%3D%20ptrI%29%20%7B%0A%20%20%20%20std%3A%3Acout%20%3C%3C%202%20%3C%3C%20'%20'%3B%0A%20%20%7D%0A%20%20if%20%28%26ptrJ1%20%3D%3D%20%26ptrJ2%29%20%7B%0A%20%20%20%20std%3A%3Acout%20%3C%3C%203%20%3C%3C%20'%20'%3B%0A%20%20%7D%0A%20%20if%20%28*ptrJ1%20%3D%3D%20*ptrI%29%20%7B%0A%20%20%20%20std%3A%3Acout%20%3C%3C%204%20%3C%3C%20'%20'%3B%0A%20%20%7D%0A%0A%20%20return%200%3B%0A%7D&cumulative=false&curInstr=0&heapPrimitives=nevernest&mode=display&origin=opt-frontend.js&py=cpp_g%2B%2B9.3.0&rawInputLstJSON=%5B%5D&textReferences=false}{Click here to see the visualization for this code.}
\end{frame}

\begin{frame}[fragile]{Assigning a value to a \textit{dereferenced} pointer.}

A pointer must have a value before you can \textit{dereference} it (follow
the pointer).\\[1em]

\begin{columns}[T]
\column{0.5\textwidth}
\begin{lstlisting}
int *x;
*x = 3;
\end{lstlisting}
  \pause
  \centering\textbf{Error! Override the value located in some unknown address.}
\column{0.5\textwidth}
  \pause
\begin{lstlisting}
int foo;
int *x;
x = &foo;
*x = 3;
\end{lstlisting}
  \pause
  \centering\textbf{This is fine. \lstinline{x} points to \lstinline{foo}.}
\end{columns}
\end{frame}

\begin{frame}[fragile]{Pointers to Pointers}\center
% \renewcommand{\arraystretch}{4}
\begin{tabular}{|>{\columncolor{codeBackground}}l l|}
\hline
&\\
\lstinline|int *x = &num;| & \raisebox{-.5\dimexpr\totalheight-\ht\strutbox}{\onslide<2->{%
\begin{tikzpicture}[pointer/.append style = {thick}, node of list/.append style = {thick, node distance=10mm}]
  \LinkPtr{x}
  \node[node of list, right = of ptr, name=firstElValue, minimum width=20mm,] {?};
  \node[inner sep=0pt, align=center, anchor=south] at ([yshift=1mm] firstElValue.north) {some \texttt{int}};
  \draw[link] (ptr) to[bend left=8] (firstElValue);
\end{tikzpicture}}}\\
&\\ \hline
&\\
\lstinline|double *y;| & \raisebox{-.5\dimexpr\totalheight-\ht\strutbox}{\onslide<3->{%
\begin{tikzpicture}[pointer/.append style = {thick}, node of list/.append style = {thick, node distance=10mm}]
  \LinkPtr{y}
  \node[node of list, right = of ptr, name=firstElValue, minimum width=30mm,] {?};
  \node[inner sep=0pt, align=center, anchor=south] at ([yshift=1mm] firstElValue.north) {some \texttt{double}};
  \draw[link] (ptr) to[bend left=8] (firstElValue);
\end{tikzpicture}}} \\
&\\ \hline
&\\
\lstinline|int **z = &x;| & \raisebox{-.5\dimexpr\totalheight-\ht\strutbox}{\onslide<4->{%
\begin{tikzpicture}[pointer/.append style = {thick}, node of list/.append style = {thick}]
  \LinkPtr{z}
  
  \node[pointer, name=ptr2, node distance=10mm, right = of ptr2] {};
  \node[node of list, node distance=10mm, right = of ptr2, name=firstElValue, minimum width=20mm,] {?};
  
  \fill (ptr2) circle(2pt);
  
  \node[inner sep=0pt, align=center, anchor=south] at ([yshift=1.5mm] ptr2.north) {\texttt{x}};

  \node[inner sep=0pt, align=center, anchor=south] at ([yshift=1mm] firstElValue.north) {some \texttt{int}};

  \draw[link] (ptr) to[bend left=8] (ptr2);
  \coordinate (ptr2) at (ptr2);
  \draw[link] (ptr2) to[bend left=8] (firstElValue);
\end{tikzpicture}}} \\
& \\ \hline
\end{tabular}

\end{frame}

\begin{frame}[fragile]{Pointers to Pointers: An Example}


\begin{lstlisting}
int num = 5;
int *pNum = &num;
int **ppNum = &pNum;
\end{lstlisting}

Then:
\begin{itemize}
\item
  \lstinline{ppNum} stores the memory location of\alt<2->{ \lstinline{pNum}.}{\textellipsis}
\item<3->
  \lstinline{*ppNum} stores the memory location of\alt<4->{ \lstinline{num}.}{\textellipsis}
\item<5->
  \lstinline{**ppNum} equals\alt<6->{ \lstinline{5}.}{\textellipsis}
\end{itemize}
\onslide<6->{}
\end{frame}

\section{Array Pointers}

\begin{frame}[fragile]{Declaring Arrays}

\begin{itemize}
\item
  Typical C/C++ array declarations.
  \begin{lstlisting}
int arr[5];               // stack
double arr1[10][15];      // stack
int *intPtr = new int[5]; // heap
\end{lstlisting}
\item
  Typical Java array declarations:
  \begin{lstlisting}
int[] arr = new int[5];
double[][] arr1 = new double [10][5];
Object[] arr2 = new Object[100];
\end{lstlisting}
\end{itemize}
\end{frame}

\begin{frame}{Allocation of Stack and Heap Space for an Array}
\begin{center}
\lstinline{int *arr = new int[10];}\\\vspace{1em}

\begin{tikzpicture}
	\ArrayLabel{}
	\begin{scope}[array index/.append style={font=\sffamily,text depth=0.2ex}, node of array/.append style={minimum width=11mm, fill=MyPurple}]
		\ArrayNode<size>{10}
		\node[block] (firstEl) [fit=(prev)]{};
		% \ArrayNode<type>{int}
		\ArrayGroup{firstEl}{firstEl}{metadata}
	\end{scope}
	\setcounter{arrayIndex}{0}
	
	\ArrayNode{?}
	\node[block] (firstArrayEl) [fit=(val)]{};
	\ArrayNode{?}
	\coordinate (heapBottom) at ([yshift=-1.5em] val.south);
	\ArrayNodeSpan{7}
	\ArrayNode{?}


	\begin{scope}[on background layer]
		\node[draw,fill=CSUslate, very thick,fit=(firstEl) (index) (groupLabel) (val) (\detokenize{el?}) (heapBottom)] (heapMembers) {};
		% \node[north=of heapMembers]{Heap};
		\node[anchor=south west, font=\it] at ([yshift=0pt] heapMembers.north west) {Heap};
		\node[anchor=east] at ([yshift=0pt] heapMembers.north west) {$h$};
		\node[anchor=west] at ([yshift=0pt] heapMembers.south east) {$n$};
	\end{scope}

	\begin{scope}[node of array/.append style={minimum width=14mm, font=\ttfamily\bfseries}]
		\ArrayLabel[-1.51][-2]{}
		%\ArrayNodeStack[CSUcyan][arr]{int}
		%\ArrayNodeStack[CSUcyan][size]{10}

		\ArrayNodeStack[CSUgold][pArr]{}
		% \ArrayGroupVerticleRight{val}{val}{Call\\Stack}
		\fill (prev) circle(3pt);
		\coordinate (ptrCircle) at (prev);
		\draw[link, shorten >=1pt] (ptrCircle) to[bend right=30] (firstArrayEl);

		\draw [densely dashed, very thick] ([xshift=0.5pt] val.north west) -- ([yshift=1.5em, xshift=0.5pt] val.north west);
		\draw [densely dashed, very thick,] ([xshift=-0.5pt] val.north east) -- ([yshift=1.5em, xshift=-0.5pt] val.north east);

		\node[anchor=south, font=\it, align=center] at ([yshift=1.5em] val.north) {Call\\Stack};
	\end{scope}
\end{tikzpicture}
\end{center}
\end{frame}

\begin{frame}{Pointers and Arrays}

\begin{itemize}
\item
  The identifier of an array on the stack is basically a \\\emph{const pointer}
  pointing at the beginning of the array.
\item
  You can use the \lstinline{{[}{]}} operator with pointers!
\item
  Example:\\
  \lstinline{int A{[}5{]};}
  \begin{itemize}
  \item
    Creates a memory block of 5 integers on the stack ($5\times4$ bytes).
  \item
    \lstinline{A} (the pointer) points at the beginning of the array. A $\rightarrow$ \lstinline{A{[}0{]}}.\\
    \lstinline{(A == \&A[0])}

    \begin{tikzpicture}
      \ArrayLabel{A}
      \ArrayList{.,.,.,.,.}
    \end{tikzpicture}
  \end{itemize}
\end{itemize}
\end{frame}

\begin{frame}[fragile]{Pointers and Arrays}
\begin{lstlisting}
int *x;
int a[5] {-1, -2, -3, -4, -5};
x = &a[2]; // x is the address of a[2]

for (int i = 0; i < 3; i++)
  x[i]++; // x[i] is the same as a[i+2]
\end{lstlisting}
\vspace{1em}
% \begin{tabularx}{0.7\textwidth}{r *{5}{>{\centering\arraybackslash}X}}
%   & \texttt{a[0]} & \texttt{a[1]} & \texttt{a[2]} & \texttt{a[3]} & \texttt{a[4]}\\
%   \cline{2-6}
%     \lstinline{a} $\rightarrow$ &
%     \multicolumn{1}{|c}{-1} &
%     \multicolumn{1}{|c}{-2} &
%     \multicolumn{1}{|c}{-3} &
%     \multicolumn{1}{|c}{-4} &
%     \multicolumn{1}{|c|}{-5} \\
%   \cline{2-6}
%   & & & $\uparrow$ & $\uparrow$ & $\uparrow$\\
%   & & & \texttt{x} \texttt{\&x{[}0{]}} & \texttt{\&x{[}1{]}} & \texttt{\&x{[}2{]}}\\
% \end{tabularx}
  \centerline{%
  \begin{tikzpicture}[node of array/.append style={minimum width=22mm}, label of array el/.style={font=\ttfamily}]
    \ArrayLabel{a}
    \ArrayList{-1, -2}
    \alt<-1>{\ArrayNode{-3}}{%
      \begin{scope}[node of array/.append style={font=\bfseries, fill=MyPurple}]%
        \ArrayNode{-2}%
      \end{scope}%
    }%
    
    \ArrayElLabel<-8mm><-6mm>{x}
    \ArrayElLabel<-8mm><6mm>{\&x{[}0{]}}
    \alt<-1>{\ArrayNode{-4}}{%
      \begin{scope}[node of array/.append style={font=\bfseries, fill=MyPurple}]%
        \ArrayNode{-3}%
      \end{scope}%
    }
    \ArrayElLabel<-8mm>{\&x{[}1{]}}
    \alt<-1>{\ArrayNode{-5}}{%
      \begin{scope}[node of array/.append style={font=\bfseries, fill=MyPurple}]%
        \ArrayNode{-4}%
      \end{scope}%
    }
    \ArrayElLabel<-8mm>{\&x{[}2{]}}
  \end{tikzpicture}%
  }%
  \onslide<2>{}
\end{frame}

\begin{frame}[fragile]{Pointer Arithmetic}
\begin{itemize}
\item
  Integer arithmetic (\texttt{+}, \texttt{-}, \texttt{++}, \texttt{--}, \texttt{+=}, \texttt{-=}) works with pointers.
\item
  Increment updates the address to the ``next'' element.
  \begin{lstlisting}
int a[5] {-1, -2, -3, -4, -5};
int *ptr = a;
*(ptr + 3) = 400;
\end{lstlisting}
\end{itemize}



% \begin{tabularx}{1\textwidth}{r *{5}{>{\centering\arraybackslash}X}}
%   & \texttt{a[0]} & \texttt{a[1]} & \texttt{a[2]} & \texttt{a[3]} & \texttt{a[4]}\\
%   \cline{2-6}
%     \lstinline{a} $\rightarrow$ &
%     \multicolumn{1}{|c}{-1} &
%     \multicolumn{1}{|c}{-2} &
%     \multicolumn{1}{|c}{-3} &
%     \multicolumn{1}{|c}{\textbf{400}} &
%     \multicolumn{1}{|c|}{-5} \\
%   \cline{2-6}
%   & $\uparrow$ & & $\uparrow$ & $\uparrow$ & $\uparrow$\\
%   & \lstinline{ptr} & & \lstinline{(ptr + 2)} & \lstinline{(ptr + 3)} & \lstinline{(ptr + 4)}\\
% \end{tabularx}

\centerline{%
  \begin{tikzpicture}[node of array/.append style={minimum width=24mm}, label of array el/.style={font=\ttfamily}]
    \ArrayLabel{a}
    \ArrayNode{-1}
    \ArrayElLabel<-7mm>{ptr}
    \ArrayNode{-2}
    \ArrayNode{-3}
    \ArrayElLabel<-7mm>{(ptr + 2)}
    \alt<-1>{\ArrayNode{-4}}{%
      \begin{scope}[node of array/.append style={font=\bfseries, fill=MyPurple}]%
        \ArrayNode{400}%
      \end{scope}%
    }
    \ArrayElLabel<-7mm>{(ptr + 3)}
    \ArrayNode{-5}
    \ArrayElLabel<-7mm>{(ptr + 4)}
  \end{tikzpicture}%
  }%
  \onslide<2>{}
\end{frame}

\section{Memory Errors}

\begin{frame}[fragile]{Pointer Pitfalls}

Assigning values to uninitialized, null, or deleted pointers:
\begin{columns}[t]
\column{0.17\textwidth}
\begin{lstlisting}[basicstyle=\ttfamily\small]
int* p;
*p = 3;
\end{lstlisting}
\column{0.35\textwidth}
\begin{lstlisting}[basicstyle=\ttfamily\small]
int* p = nullptr;
*p = 4;
\end{lstlisting}
\column{0.35\textwidth}
\begin{lstlisting}[basicstyle=\ttfamily\small]
int* p = new int;
delete p;
*p = 5;
\end{lstlisting}
\end{columns}

Using the above statements will result in a \textit{segmentation fault} or undefined behavior!

\end{frame}

\begin{frame}[fragile]{Memory Leaks}


Memory leak is when you remove the reference to the memory block,
before deleting the block itself.\vspace{1em}

Example:
\begin{lstlisting}
int *p;      // p points somewhere
p = new int; // p is an int value's address
p = nullptr; // p points to null.
\end{lstlisting}

\begin{itemize}[<+->]
\item
  Result?
\item\textbf{Memory Leek! (Orphan Blocks)}\\
  Must free memory block before changing reference.
\end{itemize}
\end{frame}


\begin{frame}{Memory Leaks}
\begin{itemize}
\item
  A memory leak can diminish the performance of the computer by
  reducing the amount of available memory.
\item
  Eventually, in the worst case, too much of the available memory
  may become allocated
\item
  and all or part of the system or devices stops working correctly,
  the application fails, or the system slows down.
\end{itemize}
\end{frame}

\begin{frame}{Memory Leaks}
  \begin{block}{Valgrind}
  Use Valgrind with \texttt{-{}-leak-check=yes} option if your implementation has
  memory leaks.\\
  Reference the \href{https://valgrind.org/docs/manual/manual.html}{Valgrind User Manual}, 
  \href{https://valgrind.org/docs/manual/quick-start.html}{The Valgrind Quick Start Guide}, or 
  a \href{https://valgrind.org/downloads/guis.html}{Graphical User Interfaces}.
  
  * This is why the \texttt{-g} flag is used when debugging!
\end{block}
\end{frame}

\begin{frame}{Problem: Multiple Pointers to the Same Address}

\begin{itemize}
\item
  A second problem can occur when multiple pointers are assigned to a block of heap memory.
\item
  The block may be deleted and one of the pointers set to \textit{null},
  but the other pointers still exist.
\item
  If the runtime system reassigns the memory to another object, the
  original pointers pose a danger.
\end{itemize}
\end{frame}

\begin{frame}[fragile]{Dangling Pointer: A pointer that points to an invalid location.}

Example:

\begin{lstlisting}
int *p, *q;  // Create two pointers
p = new int; // Allocate an int on the heap.
q = p;       // Points to the same address.
*p = 4;      // Sets the on the heap to 4.
delete q;    // Frees the heap memory.
*p = 3;      // Illegal assignment!
\end{lstlisting}

\begin{itemize}
\item
  \lstinline{p} and \lstinline{q} point to the same location.
\item
  When \lstinline{q} is deleted, \lstinline{p} becomes a dangling pointer!
\end{itemize}
\end{frame}

\section{Garbage Collection}
\begin{frame}{Memory Leaks and Garbage Collection}
  The popularity of \textit{object-oriented programming} has meant more emphasis
  on heap storage management.\\[1em]

  Objects are considered active or inactive.
  \begin{itemize}
  \item
    \emph{Active} objects: blocks \textit{accessible} through a pointer or reference.
  \item
    \emph{Inactive} objects: \textit{inaccessible} blocks; no reference exists (also called (\emph{orphans} or \emph{garbage}).
  \end{itemize}

  The terms \textit{accessible} and \textit{inaccessible} may be more figurative than objectively true.
\end{frame}

\begin{frame}{Garbage Collection}

\begin{itemize}
\item
  All inaccessible blocks of storage are identified and returned to the
  free list.
\item
  The heap may also be \textbf{compacted} at this time: allocated space is
  compressed into one end of the heap, leaving all free space in a large
  block at the other end.
\end{itemize}
\end{frame}

\begin{frame}{Garbage Collection}

\begin{itemize}[<+->]
\item
  C \& C++ leave it to the programmer --- if an unused storage block
  isn't explicitly freed by, it becomes garbage.
  \begin{itemize}
  \item
    You can get C++ garbage collectors, but they aren't standard.
  \end{itemize}
\item
  Java, Python, Perl, and other scripting languages perform garbage collection.
  \begin{itemize}
  \item
    Python, etc.\ also have automatic allocation, so the ``new'' operator is not needed).
  \end{itemize}
\item
  Garbage collection was pioneered by languages like
  \href{https://en.wikipedia.org/wiki/Lisp_(programming_language)}{Lisp}, which
  constantly creates and destroys linked lists.

\item C++ includes \href{https://en.cppreference.com/w/cpp/memory}{smart pointers}, enabling automatic object lifetime management.
\end{itemize}
\end{frame}


\section{Review}

\begin{frame}{Review}

\begin{itemize}[<+->]
\item
  Three types of memory storage:
  \begin{itemize}
  \item
    Static
  \item
    Stack
  \item
    Heap
  \end{itemize}
\item
  Problems with heap storage:
  \begin{itemize}
  \item
    Memory leaks (garbage): failure to free storage when pointers
    (references) are reassigned.
  \item
    Dangling pointers: when storage is freed, but references to the
    storage still exist.
  \end{itemize}
\end{itemize}
\end{frame}

\begin{frame}{Allocation of Stack and Heap Space for an Array}
\begin{center}
  \lstinline{int *arr = new int[10];}\\\vspace{1em}

  % \includegraphics[width=0.8\textwidth]{memory-allocation-heap}%
  \begin{tikzpicture}
	\ArrayLabel{}
	\begin{scope}[array index/.append style={font=\sffamily,text depth=0.2ex}, node of array/.append style={minimum width=11mm, fill=MyPurple}]
		\ArrayNode<size>{10}
		\node[block] (firstEl) [fit=(prev)]{};
		% \ArrayNode<type>{int}
		\ArrayGroup{firstEl}{firstEl}{metadata}
	\end{scope}
	\setcounter{arrayIndex}{0}
	
	\ArrayNode{?}
	\node[block] (firstArrayEl) [fit=(val)]{};
	\ArrayNode{?}
	\coordinate (heapBottom) at ([yshift=-1.5em] val.south);
	\ArrayNodeSpan{7}
	\ArrayNode{?}


	\begin{scope}[on background layer]
		\node[draw,fill=CSUslate, very thick,fit=(firstEl) (index) (groupLabel) (val) (\detokenize{el?}) (heapBottom)] (heapMembers) {};
		% \node[north=of heapMembers]{Heap};
		\node[anchor=south west, font=\it] at ([yshift=0pt] heapMembers.north west) {Heap};
		\node[anchor=east] at ([yshift=0pt] heapMembers.north west) {$h$};
		\node[anchor=west] at ([yshift=0pt] heapMembers.south east) {$n$};
	\end{scope}

	\begin{scope}[node of array/.append style={minimum width=14mm, font=\ttfamily\bfseries}]
		\ArrayLabel[-1.51][-2]{}
		%\ArrayNodeStack[CSUcyan][arr]{int}
		%\ArrayNodeStack[CSUcyan][size]{10}

		\ArrayNodeStack[CSUgold][pArr]{}
		% \ArrayGroupVerticleRight{val}{val}{Call\\Stack}
		\fill (prev) circle(3pt);
		\coordinate (ptrCircle) at (prev);
		\draw[link, shorten >=1pt] (ptrCircle) to[bend right=30] (firstArrayEl);

		\draw [densely dashed, very thick] ([xshift=0.5pt] val.north west) -- ([yshift=1.5em, xshift=0.5pt] val.north west);
		\draw [densely dashed, very thick,] ([xshift=-0.5pt] val.north east) -- ([yshift=1.5em, xshift=-0.5pt] val.north east);

		\node[anchor=south, font=\it, align=center] at ([yshift=1.5em] val.north) {Call\\Stack};
	\end{scope}
\end{tikzpicture}%
\end{center}
\end{frame}

\begin{frame}{Garbage}

\begin{itemize}
\item
  \emph{Garbage}: any block of heap memory that cannot be accessed by the
   program (i.e., there is no stack pointer to the block) but in which the 
   runtime system thinks is in use.
\item
  Garbage is created in several ways:
  \begin{itemize}
  \item
    A function ends without returning the space allocated to a local array
    or other dynamic variable. The pointer is gone.
  \item
    A node is deleted from a linked data structure, but isn't freed.
  \item
    \textellipsis{}
  \end{itemize}
\end{itemize}
\end{frame}



\begin{frame}{Terminology}

\begin{itemize}
\item
  A \textbf{dangling pointer} (or dangling reference, or widow) is a pointer
  (reference) that contains the address of previously deallocated/freed heap space.
\item
  An \textbf{orphan (or garbage)} is a block of allocated heap memory that is no
  longer accessible through any pointer.
\item
  A \textbf{memory leak} is a gradual loss of available memory due to the
  creation of garbage.
\end{itemize}
\end{frame}



\section*{Lab 3}

\begin{frame}{Lab 3: Pointers}
\begin{center}\Large
Let's take a look at Lab 3 (on Blackboard).
\end{center}
\end{frame}

\end{document}